% Teil II: der durchgerechnete Weg.

% Gedankenschritt: Ungleichung nach n auflösen; Teilfolgen zeigen
% Nichtkonvergenz; geometrische Reihe durch Index-Korrektur.

\subsection*{Gegeben und gesucht}
(a) $a_n=\frac1{\sqrt{n^3+2}}$: konkretes $N$ f\"ur $|a_n|<10^{-3}$ und ein
von $\varepsilon$ abh\"angiges $N$; (b) warum $b_n=(-1)^n$ nicht gegen $1$
oder $-1$ konvergiert; (c) $\sum_{k=2}^{\infty}(-1)^k\frac{3^k}{4^{k+1}}$.

\subsection*{a) konkretes $N$ f\"ur $|a_n|<10^{-3}$}
Es ist $a_n>0$. Gefordert $\frac1{\sqrt{n^3+2}}<10^{-3}$:
\begin{equation}
  \sqrt{n^3+2}>1000\ \Longleftrightarrow\ n^3+2>10^6 .
\end{equation}
F\"ur $n\ge100$ gilt $n^3+2\ge100^3+2=10^6+2>10^6$. W\"ahle $N=99$; f\"ur
alle $n>N$ (d.h.\ $n\ge100$) ist $|a_n|<10^{-3}$.

\subsection*{a') allgemeines $N$ f\"ur $|a_n|<\varepsilon$}
\begin{equation}
  \frac1{\sqrt{n^3+2}}<\varepsilon\ \Longleftrightarrow\
  n^3+2>\frac1{\varepsilon^2}\ \Longleftrightarrow\
  n>\sqrt[3]{\frac1{\varepsilon^2}-2}.
\end{equation}
W\"ahle
\begin{equation}
  N=\left\lceil\sqrt[3]{\frac1{\varepsilon^2}-2}\right\rceil .
\end{equation}
F\"ur $\varepsilon=10^{-3}$: $N=\lceil\sqrt[3]{10^6-2}\rceil
=\lceil99.9999\rceil=100$ --- und mit der Bedingung ``$n>N$`` nimmt man
$N=99$, sodass $n\ge100$ gilt. \textbf{(Die Wahl ``$n>N$`` oder ``$n\ge N$``
verschiebt $N$ um $1$, das Ergebnis bleibt korrekt.)}

\subsection*{b) $b_n=(-1)^n$}
Teilfolgen: $b_{2k}=(-1)^{2k}=1$ und $b_{2k+1}=(-1)^{2k+1}=-1$. W\"are
$b_n$ konvergent, m\"ussten beide Teilfolgen gegen \emph{denselben}
Grenzwert gehen; da $1\neq-1$, ein Widerspruch. Also konvergiert $b_n$
weder gegen $1$ noch gegen $-1$.

\subsection*{c) $\sum_{k=2}^{\infty}(-1)^k\frac{3^k}{4^{k+1}}$}
$4^{k+1}=4\cdot4^k$, also
\begin{equation}
  \sum_{k=2}^{\infty}(-1)^k\frac{3^k}{4^{k+1}}
  =\frac14\sum_{k=2}^{\infty}\Bigl(-\frac34\Bigr)^k .
\end{equation}
Geometrische Summe ab $k=0$: $\sum_{k=0}^{\infty}(-\tfrac34)^k
=\frac1{1+\frac34}=\frac47$; die Glieder $k=0,1$ sind $1$ und
$-\frac34$:
\begin{equation}
  \sum_{k=2}^{\infty}\Bigl(-\frac34\Bigr)^k
  =\frac47-1+\frac34=\frac9{28},
  \qquad
  \frac14\cdot\frac9{28}=\frac9{112}.
\end{equation}

\begin{ergebnis}
\begin{equation*}
  \text{a)}\ N=99\ (\varepsilon=10^{-3}),\quad
  N=\left\lceil\sqrt[3]{\tfrac1{\varepsilon^2}-2}\right\rceil,\qquad
  \text{b)}\ b_n\ \text{divergent},\qquad
  \text{c)}\ \frac9{112}.
\end{equation*}
\end{ergebnis}

\begin{kontrolle}
\textbf{a) Randpunkte numerisch:} bei $n=99$:
$|a_{99}|=\frac1{\sqrt{970299+2}}>10^{-3}$ (Fenster noch nicht
erreicht); bei $n=100$: $|a_{100}|<10^{-3}$. Passt exakt zur Wahl
$N=99$, $n\ge100$.

\textbf{c) Gr\"o\ss{}enordnung:} Das erste Glied ($k=2$) ist
$+\frac{9}{64}\approx0.14$, das Vorzeichenwechselt weiter; die Summe
mindestens $\frac{9}{112}\approx0.080$. Grobe Schranke:
$\frac14\cdot\frac{(3/4)^2}{1-3/4}=\frac14\cdot\frac{9/16}{1/4}
=\frac9{16}$ (Betragsreihe) --- der tats\"achliche Wert $0.080$ liegt
darunter, weil alternierend gek\"urzt. Plausibel.

\textbf{Grenzfall / Dezimal:} $9/112=0.0803571\dots$; numerisch summiert
man die ersten 10000 Glieder: ~1/4 * 0.3214? = about 0.0804, stimmt mit
$\frac9{112}$ \"uberein.
\end{kontrolle}

\begin{merksatz}
Konvergenz-Beweise f\"uhren das $N$ auf die $\varepsilon$-Forderung
zur\"uck (Ungleichung nach $n$ aufl\"osen); alternierende Folgen widerlegt
man \"uber Teilfolgen; gebrochene geometrische Reihen klammert man und
korrigiert den Startindex.
\end{merksatz}
